
The following definitions are presented from standard mathematical sources (\textbf{Sets}).
Clarifications and labels have been added where applicable.

\begin{define}{Basic Set Concepts}
  \label{def:basic-set-concepts}
  \begin{enumerate}
    \item \textbf{Universal Set:} A universal set $\setUNIV$ is a
    set that contains all objects under consideration in a particular context or problem.
    \item \textbf{Empty Set:} The empty set, denoted $\setNULL$,
    is the unique set that contains no elements.
    \item \textbf{Subset:} Let $A$ and $B$ be sets. $A$ is a
    subset of $B$ ($A \subseteq B$) if (and only if)
    for all $x \in A$, $x \in B$.
    \item \textbf{Set Equality:} Let $A$ and $B$ be sets. The sets $A$ and $B$ are
    equal ($A = B$) if $A \subseteq B$ and $B \subseteq A$.
    \item \textbf{Complementary Set:} Let $A$ be a subset of the universal set $\setUNIV$.
    The complement of $A$, denoted $A^c$ or $\overline{A}$, is the set of all elements
    in $\setUNIV$ that are not in $A$. That is, $A^c = \{x \in \setUNIV : x \notin A\}$.
  \end{enumerate}
\end{define}

The following definitions extend the understanding of set operations
and properties.

\begin{define}{Set Operations}
  \label{def:set-operations}
  \begin{enumerate}
    \item \textbf{Set Difference:} Let $A$, $B$ be subsets of a universal set $\setUNIV$.
    The set difference $A \sMINUS B$ is defined to be the set
    $\{ x \in \setUNIV \LDIV x \in A \text{ and } x \notin B\}$. \vspace*{10px}\\
    \textit{It is sometimes called the relative complement of $B$ in $A$,
    and denoted in the textbook as $A - B$.}
    \item \textbf{Cartesian Product:} Let $A$ and $B$ be sets. The Cartesian product of
    $A$ and $B$, denoted $A \sTIMES B$, is the set consisting of all ordered pairs:
    \[
    A \sTIMES B = \{(a, b) \LDIV a \in A, b \in B\}.
    \]
  \end{enumerate}
\end{define}


\newpage

The following claims are given in the reference as well. Let $A,B\in U$, a universal set.

\begin{claim}{Subset of Union: $A \subseteq \sOR{A}{B}$}
  \label{claim:subset-union}
  \begin{subprove}[bluecolor]
      Let $x\in U$, assume $x\in A$. Then the statement \inlinequote{$x\in A$ or $x\in B$} is true.
      Hence, by the definition, $\sOR{A}{B}=\{x\in U\LDIV x\in A$ or $x\in B\}$, $x\in \sOR{A}{B}$.
      Thus, $A \subseteq \sOR{A}{B}$.
  \end{subprove}
\end{claim}

\begin{claim}{Empty Subset: $\setNULL \subseteq A$}
  \label{claim:empty-subset}
  \begin{subprove}[bluecolor]
      Let $x\in U$, assume $x\in \setNULL$. But there are no elements in $\setNULL$, so
      $x\in \setNULL$ is false. Since the premise is false, the implication if $x\in \setNULL$,
      then $x\in A$ is true. Thus, $\setNULL \subseteq A$.
  \end{subprove}
\end{claim}

\begin{claim}{Subset as Equality: $\sOR{A}{B}=A$ if and only if $B\subseteq A$}
  \label{claim:subset-equality}
  \begin{subprove}[bluecolor]
      Let $x\in U$, first, we assume $B\subseteq A$, we have already shown
      $A \subseteq \sOR{A}{B}$ by Claim~\ref{claim:subset-union}. To show $\sOR{A}{B} \subseteq A$, assume $x\in \sOR{A}{B}$.
      By definition, $x\in A$ or $x\in B$. Considering $x\in B$ and assumption
      $B \subseteq A$, $x\in B$ implies $x\in A$. Thus, $\sOR{A}{B} \subseteq A$.
      Hence, if $B\subseteq A$, then $\sOR{A}{B} = A$. \vspace*{10px}\\
      Now, we assume $\sOR{A}{B}=A$. To show $B\subseteq A$, we now assume $x\in B$. Since $x\in B$, by definition
      of union, $x\in \sOR{A}{B} = A$. Therefore, $B\subseteq A$.
  \end{subprove}
\end{claim}